\begin{keybox}
\textbf{The answer in short.} You asked whether a model that is kept the same size does worse
when its language has more words. I trained small models on a made-up language about animals,
with 16 to 1,024 animal words. The models came in four sizes, which I call widths 4, 8, 16 and
32. Each was graded on how many facts about the animals it knew, on a score from 0 (no better
than guessing) to 100 (every fact right).

\textbf{The verdict by the rule we fixed in advance: undecided.} The rule asks for a model
that learns the small language, comes to a stop clearly lower on a larger one, and is beaten
there by a model twice as wide. No model did that. One came close (point 5), but by the
test we fixed for it none of its runs had stopped improving when their time ran out. (The rule could only ever come out in favor of your idea or undecided.)

\begin{enumerate}[itemsep=1pt,topsep=2pt]
\item \textbf{The planned experiment did not put your idea to the test} (section
      \ref{sec:round1}). Widths 8, 16 and 32 each learned every fact of every language, from 16
      animals to 1,024, in all 36 runs. By a rough guide none of them was much more than
      one-third full, so there was nothing for more words to push out. What more words did
      cost was training: 64 times as many animals took about 6 to 8 times as much text.
\item \textbf{Using some animals far more often than others did no harm} (section
      \ref{sec:zipf}). All 36 runs were solved again, and the large languages were learned
      sooner than with equal use.
\item \textbf{Part-way through training, the scores look like your prediction} (section
      \ref{sec:equal}). This is the comparison you suggested: every model after the same
      amount of text. At one such point width 16 scored 100 on 16 animals but only 59 on 256
      animals, and width 32 scored 81 on those 256. But the cause is not the one your idea
      names. Your idea is about room; this is about speed. A larger language takes more text
      to learn and a wider model learns it from less, and with more text every one of these
      models reached 100.
\item \textbf{The claim your idea starts from, that a model stores what words mean, is
      supported} (section \ref{sec:probe}). Each word has a row of numbers of its own in the
      model. A simple test could read an animal's facts from its row 70 to 97 times in 100,
      and animals with the same facts ended up with nearly the same row.
\item \textbf{The nearest thing to your prediction came when each animal had 16 facts in
      place of 4} (section \ref{sec:round32}). Widths 8, 16 and 32 all still learned the 16
      animals. On 256 animals widths 16 and 32 learned practically everything too, but width 8
      did not. In three runs it stayed below 60, and in the two that got furthest it held only
      about two of the four kinds of fact. That is what a model too small for its language would
      look like. One night was not enough to tell it apart from a model that is only very slow.
\item \textbf{Runs can stand still for a long time and then go on} (section
      \ref{sec:round32}). In that same round two models stood at a score of 37 for so long
      that our test for the end of a run called them finished. I replayed one of them exactly
      and let it run four times as long: it stood at 37 for most of that time and then rose
      to 98.6. The other I started again from the same numbers on the same text, with the
      computer set up slightly differently, which changes only the last digits of its
      arithmetic. It never paused at 37 and reached 100. Whether a run got stuck was chance.
\item \textbf{The smallest model, width 4, learned no facts in the time it had} (section
      \ref{sec:four}), not even about 16 animals, though it learned the shapes of the
      sentences. It seems too narrow to work at all, which would be a different thing from
      being too full.
\end{enumerate}

\textbf{Two lessons.} First, because every word has a row of its own, each added word brings
room with it. That is the likely reason why adding words alone did not fill a model here. To
test your prediction properly, the words must be made to share one store of fixed size;
section \ref{sec:extras} says how. Second, these models learn in jumps, with long stand-stills
in between. A model that has stopped improving is not necessarily finished.
\end{keybox}
